%=============================================================================
% Chapter 1 Exercises
%=============================================================================
\section{Chapter 1 Exercises [240]}

%-----------------------------------------------------------------------------
\subsection{[2] Types of Computers}
\hypertarget{prob-1.1}{}
\noindent$\langle$\S1.1$\rangle$\medskip

\noindent List and describe three types of computers.

\problembreak % next [5] → 2+5=7 ≥ 5

%-----------------------------------------------------------------------------
\subsection{[5] Seven Great Ideas Matching}
\hypertarget{prob-1.2}{}
\noindent$\langle$\S1.2$\rangle$\medskip

\noindent The seven great ideas in computer architecture are similar to ideas from other
fields. Match the seven ideas from computer architecture, ``Use Abstraction to
Simplify Design'', ``Make the Common Case Fast'', ``Performance via
Parallelism'', ``Performance via Pipelining'', ``Performance via Prediction'',
``Hierarchy of Memories'', and ``Dependability via Redundancy'' to the
following ideas from other fields:

\begin{enumerate}
  \item[a.] Assembly lines in automobile manufacturing
  \item[b.] Suspension bridge cables
  \item[c.] Aircraft and marine navigation systems that incorporate wind
        information
  \item[d.] Express elevators in buildings
  \item[e.] Library reserve desk
  \item[f.] Increasing the gate area on a CMOS transistor to decrease its
        switching time
  \item[g.] Self-driving cars whose control systems primarily rely on
        existing sensor systems already installed into the base vehicle,
        such as heated destination systems with driver and front-seat controls
\end{enumerate}

\problembreak % [5] ≥ 5

%-----------------------------------------------------------------------------
\subsection{[2] High-Level Language to Machine Code}
\hypertarget{prob-1.3}{}
\noindent$\langle$\S1.3$\rangle$\medskip

\noindent Describe the steps that transform a program written in a high-level language
such as C into a representation that is directly executed by a computer
processor.

\ifprintversion\vfill\else\bigskip\fi % grouped with 1.4: 2+2=4 < 5

%-----------------------------------------------------------------------------
\subsection{[2] Color Display Frame Buffer}
\hypertarget{prob-1.4}{}
\noindent$\langle$\S1.4$\rangle$\medskip

\noindent Assume a color display using 8 bits for each of the primary colors (red,
green, blue) per pixel and a frame size of $640 \times 480$.

\begin{enumerate}
  \item[a.] What is the minimum size in bytes of the frame buffer to store a
        frame?
  \item[b.] How long would it take, at a minimum, for the frame to be sent
        over a 100 Mbit/s network?
\end{enumerate}

\problembreak % next [4] → 4+4=8 ≥ 5

%-----------------------------------------------------------------------------
\subsection{[4] Comparing Three Processors}
\hypertarget{prob-1.5}{}
\noindent$\langle$\S1.6$\rangle$\medskip

\noindent Consider three different processors P1, P2, and P3 executing the same
instruction set. P1 has a 3 GHz clock rate and a CPI of 1.5. P2 has a 2.5
GHz clock rate and a CPI of 1.0. P3 has a 4.0 GHz clock rate and a CPI of
2.2.

\begin{enumerate}
  \item[a.] Which processor has the highest performance expressed in
        instructions per second?
  \item[b.] If the processors each execute a program in 10 seconds, find the
        number of cycles and the number of instructions.
  \item[c.] We are trying to reduce the execution time by 30\% but this leads
        to an increase of 20\% in the CPI. What clock rate should we have to
        get this time reduction?
\end{enumerate}

\problembreak % next [5] → 4+5=9 ≥ 5

%-----------------------------------------------------------------------------
\subsection{[5] Intel Processor Performance Trends}
\hypertarget{prob-1.6}{}
\medskip

\noindent Consider the table given next, which tracks several performance indicators for
Intel desktop processors since 2010. The ``Tech'' column shows the minimum
feature size of each processor's fabrication process. Assume that the die size
has remained relatively constant and the number of transistors that comprise
each processor scales at $(1/T)^2$, where $T$ = the minimum feature size. For
each performance indicator, calculate the average rate of improvement from
2010 to 2019, as well as the number of years required to double each at that
corresponding rate.

\begin{center}
\begin{tabular}{lccccccc}
\toprule
\textbf{Desktop} & \textbf{Year} & \textbf{Tech} & \textbf{Clock} &
\textbf{Integer} & \textbf{Cores} & \textbf{Max.\ DRAM} &
\textbf{DRAM BW} \\
\textbf{Processor} & & \textbf{(nm)} & \textbf{(GHz)} &
\textbf{IPC/Core} & & \textbf{Capacity} & \textbf{(GB/s)} \\
\midrule
Core i7 870      & 2010 & 45  & 2.93 & 1.0 & 4  & 16 GB & 21  \\
Core i7 4770     & 2013 & 22  & 3.40 & 1.3 & 4  & 32 GB & 25.6 \\
Core i7 6700K    & 2015 & 14  & 4.00 & 1.5 & 4  & 64 GB & 34.1 \\
Core i7 8700K    & 2017 & 14  & 3.70 & 1.5 & 6  & 64 GB & 38.4 \\
Core i9 9900K    & 2019 & 14  & 3.60 & 1.5 & 8  & 128 GB & 41.6 \\
\midrule
\textbf{Imp/year} & & & & & & & \\
\bottomrule
\end{tabular}
\end{center}

\problembreak % [5] ≥ 5

%-----------------------------------------------------------------------------
\subsection{[20] Two ISA Implementations Compared}
\hypertarget{prob-1.7}{}
\noindent$\langle$\S1.6$\rangle$\medskip

\noindent Consider two different implementations of the same instruction set
architecture. The instructions can be divided into four classes according to
their CPI (columns A, B, C, and D). P1 has a clock rate of 2.5 GHz and CPIs
of 1, 2, 3, and 3. P2 with a clock rate of 3 GHz and CPI of 2, 2, 2, and 2.

Given a program with a dynamic instruction count of $1.0 \times 10^6$
instructions divided into classes as follows: 10\% class A, 20\% class B,
50\% class C, and 20\% class D, which is faster: P1 or P2?

\begin{enumerate}
  \item[a.] Which is faster: P1 or P2?
  \item[b.] Find the clock cycles required in both cases.
\end{enumerate}

\problembreak % [20] ≥ 5

%-----------------------------------------------------------------------------
\subsection{[5] Compiler Impact on Performance}
\hypertarget{prob-1.8}{}
\noindent$\langle$\S1.6$\rangle$\medskip

\noindent Compilers can have a profound impact on the performance of an application.
Assume that for a program, compiler A results in a dynamic instruction count
of $1.0 \times 10^9$ that has an execution time of 1.1 s, while compiler B
results in a dynamic instruction count of $1.2 \times 10^9$ and an execution
time of 1.5 s.

\begin{enumerate}
  \item[a.] Find the average CPI for each program given that the processor has
        a clock cycle time of 1 ns.
  \item[b.] Assume the compiled programs run on two different processors. If
        the execution times on the two processors are the same, how much is
        the clock of the processor running compiler A's code versus the clock
        of the processor running compiler B's code?
  \item[c.] A new compiler is developed that uses only $6.0 \times 10^8$
        instructions and has an average CPI of 1.1. What is the speedup of
        using this new compiler versus using compilers A or B on the original
        processor?
\end{enumerate}

\problembreak % [5] ≥ 5

%-----------------------------------------------------------------------------
\subsection{[15] Processor Power and Capacitance}
\hypertarget{prob-1.9}{}
\noindent$\langle$\S1.7$\rangle$\medskip

\noindent The Pentium 4 Prescott processor, released in 2004, had a clock rate of 3.6
GHz and voltage of 1.25 V. Assume that, on average, it consumed 10 W of
static power and 90 W of dynamic power.

The Core i5 Ivy Bridge, released in 2012, had a clock rate of 3.4 GHz and
voltage of 0.9 V. Assume that, on average, it consumed 30 W of static power
and 40 W of dynamic power.

\begin{enumerate}
  \phantomsection\addcontentsline{toc}{subsubsection}{[5] 1.9.1}
  \item[{[5] 1.9.1}] For each processor find the average capacitance.
  \phantomsection\addcontentsline{toc}{subsubsection}{[5] 1.9.2}
  \item[{[5] 1.9.2}] Find the percentage of the total dissipated power comprised by
        static power and the ratio of static power to dynamic power for each
        technology.
  \phantomsection\addcontentsline{toc}{subsubsection}{[5] 1.9.3}
  \item[{[5] 1.9.3}] If the total dissipated power is to be reduced by 10\%, how
        much should the voltage be reduced to maintain the same leakage
        current? Note: power is defined as the product of voltage and current.
\end{enumerate}

\problembreak % [5] ≥ 5

%-----------------------------------------------------------------------------
\subsection{[20] Multiprocessor Execution Speedup}
\hypertarget{prob-1.10}{}
\noindent$\langle$\S1.10$\rangle$\medskip

\noindent Assume arithmetic, load/store, and branch instructions; a processor has CPIs
of 1, 12, and 5, respectively. Also assume that on a single processor a
program requires the execution of $2.56 \times 10^9$ arithmetic instructions,
$1.28 \times 10^9$ load/store instructions, and 256 million branch
instructions.

\begin{enumerate}
  \phantomsection\addcontentsline{toc}{subsubsection}{[10] 1.10.1}
  \item[{[10] 1.10.1}] Find the total execution time for this program on a 1, 2, and
        8 processor system. Also find the relative speedup of the 2 and 8
        processor system over the single processor.
  \phantomsection\addcontentsline{toc}{subsubsection}{[10] 1.10.2}
  \item[{[10] 1.10.2}] If the time for the arithmetic instructions was reduced by
        half, what is the overall speedup and the percentage of execution time
        due to arithmetic instructions?
\end{enumerate}

\problembreak % [10] ≥ 5

%-----------------------------------------------------------------------------
\subsection{[20] Wafer Yield and Cost per Die}
\hypertarget{prob-1.11}{}
\noindent$\langle$\S1.5$\rangle$\medskip

\noindent A 15 cm diameter wafer has a cost of \$12, contains 84 dies, and has 0.020
defects/cm$^2$. Assume a 20 cm diameter wafer has a cost of \$15, contains
100 dies, and has 0.031 defects/cm$^2$.

\begin{enumerate}
  \phantomsection\addcontentsline{toc}{subsubsection}{[5] 1.11.1}
  \item[{[5] 1.11.1}] Find the yield for both wafers.
  \phantomsection\addcontentsline{toc}{subsubsection}{[5] 1.11.2}
  \item[{[5] 1.11.2}] Find the cost per die for both wafers.
  \phantomsection\addcontentsline{toc}{subsubsection}{[5] 1.11.3}
  \item[{[5] 1.11.3}] If the number of dies per wafer is increased by 10\% and the
        defects per area unit increases by 15\%, find the die area.
  \phantomsection\addcontentsline{toc}{subsubsection}{[5] 1.11.4}
  \item[{[5] 1.11.4}] Assume a fabrication process improves the yield from 0.92 to
        0.95. Find the defect density for each yield.
\end{enumerate}

\problembreak % [5] ≥ 5

%-----------------------------------------------------------------------------
\subsection{[80] SPEC CPU2006 Benchmark Analysis}
\hypertarget{prob-1.12}{}
\noindent$\langle$\S1.6$\rangle$\medskip

\noindent The results of the SPEC CPU2006 bzip2 benchmark running on an AMD Barcelona
has an instruction count of $2.389 \times 10^{12}$, an execution time of
750 s, and a reference time of 9650 s.

\begin{enumerate}
  \phantomsection\addcontentsline{toc}{subsubsection}{[5] 1.12.1}
  \item[{[5] 1.12.1}] Find the CPI if the clock cycle time is 0.333 ns.
  \phantomsection\addcontentsline{toc}{subsubsection}{[5] 1.12.2}
  \item[{[5] 1.12.2}] Find the SPECratio.
  \phantomsection\addcontentsline{toc}{subsubsection}{[5] 1.12.3}
  \item[{[5] 1.12.3}] Find the increase in CPU time if the number of
        instructions of the benchmark is increased by 10\% without
        affecting the CPI.
  \phantomsection\addcontentsline{toc}{subsubsection}{[5] 1.12.4}
  \item[{[5] 1.12.4}] Find the increase in CPU time if the number of
        instructions of the benchmark is increased by 10\% and the CPI
        is increased by 5\%.
  \phantomsection\addcontentsline{toc}{subsubsection}{[5] 1.12.5}
  \item[{[5] 1.12.5}] Find the change in the SPECratio for this change.
  \phantomsection\addcontentsline{toc}{subsubsection}{[10] 1.12.6}
  \item[{[10] 1.12.6}] Suppose that we are developing a new version of the
        AMD Barcelona processor with a 4 GHz clock rate. We have added some
        additional instructions to the instruction set in such a way that the
        number of instructions has been reduced by 15\%. The execution time
        is reduced to 700 s and the new SPECratio is 13.7. Find the new CPI.
  \phantomsection\addcontentsline{toc}{subsubsection}{[10] 1.12.7}
  \item[{[10] 1.12.7}] This CPI value is larger than obtained in 1.12.1 as
        the clock rate was increased from 3 GHz to 4 GHz. Determine whether
        the increase in the CPI is similar to that of the clock rate. If
        they are dissimilar, why?
  \phantomsection\addcontentsline{toc}{subsubsection}{[5] 1.12.8}
  \item[{[5] 1.12.8}] By how much has the CPU time been reduced?
  \phantomsection\addcontentsline{toc}{subsubsection}{[10] 1.12.9}
  \item[{[10] 1.12.9}] For a second benchmark, libquantum, assume an
        execution time of 960 s, CPI of 1.61, and clock rate of 3 GHz.
        If the execution time is reduced by an additional 10\% without
        affecting the CPI and with a clock rate of 4 GHz, determine the
        number of instructions.
  \phantomsection\addcontentsline{toc}{subsubsection}{[10] 1.12.10}
  \item[{[10] 1.12.10}] Determine the clock rate required to give a further
        10\% reduction in CPU time while maintaining the number of
        instructions and with the CPI unchanged.
  \phantomsection\addcontentsline{toc}{subsubsection}{[10] 1.12.11}
  \item[{[10] 1.12.11}] Determine the clock rate if the CPI is reduced by
        15\% and the CPU time by 20\% while the number of instructions is
        unchanged.
\end{enumerate}

\problembreak % [5] ≥ 5

%-----------------------------------------------------------------------------
\subsection{[25] Performance Equation Pitfall}
\hypertarget{prob-1.13}{}
\noindent$\langle$\S1.11$\rangle$\medskip

\noindent Section 1.11 cites as a pitfall the utilization of a subset of the
performance equation as a performance metric. To illustrate this, consider the
following: P1 has a clock rate of 4 GHz, average CPI of 0.9, and requires the
execution of $5.0 \times 10^9$ instructions. P2 has a clock rate of 3 GHz, an
average CPI of 0.75, and requires the execution of $1.0 \times 10^9$
instructions.

\begin{enumerate}
  \phantomsection\addcontentsline{toc}{subsubsection}{[5] 1.13.1}
  \item[{[5] 1.13.1}] One useful fallacy is to consider that the computer with the
        largest clock rate also has the largest performance. Check if this is
        true for P1 and P2.
  \phantomsection\addcontentsline{toc}{subsubsection}{[10] 1.13.2}
  \item[{[10] 1.13.2}] Another fallacy is to consider that the processor executing
        the largest number of instructions will need a larger CPU time.
        Considering that processor P1 is executing a sequence of
        $1.0 \times 10^9$ instructions and that the CPI of processors P1 and
        P2 do not change, determine the number of instructions that P2 can
        execute in the same time that P1 needs to execute
        $1.0 \times 10^9$ instructions.
  \phantomsection\addcontentsline{toc}{subsubsection}{[10] 1.13.3}
  \item[{[10] 1.13.3}] A common fallacy is to use MIPS (millions of instructions per
        second) to compare the performance of two different processors, and
        consider the processor with the largest MIPS as having the largest
        performance.
\end{enumerate}

\problembreak % [5] ≥ 5

%-----------------------------------------------------------------------------
\subsection{[15] Amdahl's Law and Partial Improvement}
\hypertarget{prob-1.14}{}
\noindent$\langle$\S1.11$\rangle$\medskip

\noindent Another pitfall cited in Section 1.11 is expecting to improve the overall
performance of a computer by improving only one aspect of the computer.
Consider a computer running a program that requires 250 s, with 70 s spent
executing FP instructions, 85 s executed L/S instructions, and 40 s spent
executing branch instructions.

\begin{enumerate}
  \phantomsection\addcontentsline{toc}{subsubsection}{[5] 1.14.1}
  \item[{[5] 1.14.1}] By how much is the total time reduced if the time for FP
        operations is reduced by 20\%?
  \phantomsection\addcontentsline{toc}{subsubsection}{[5] 1.14.2}
  \item[{[5] 1.14.2}] By how much is the time for INT operations reduced if the
        total time is reduced by 20\%?
  \phantomsection\addcontentsline{toc}{subsubsection}{[5] 1.14.3}
  \item[{[5] 1.14.3}] Can the total time be reduced by 20\% by reducing only the
        time for branch instructions?
\end{enumerate}

\problembreak % [5] ≥ 5

%-----------------------------------------------------------------------------
\subsection{[15] CPI Optimization for Speedup}
\hypertarget{prob-1.15}{}
\noindent$\langle$\S1.6$\rangle$\medskip

\noindent A program requires the execution of $50 \times 10^6$ FP instructions,
$110 \times 10^6$ INT instructions, $80 \times 10^6$ L/S instructions, and
$16 \times 10^6$ branch instructions. The CPI for each type of instruction is
1, 1, 4, and 2, respectively. Assume that the processor has a 2 GHz clock.

\begin{enumerate}
  \phantomsection\addcontentsline{toc}{subsubsection}{[5] 1.15.1}
  \item[{[5] 1.15.1}] By how much must we improve the CPI of FP instructions if we
        want the program to run two times faster?
  \phantomsection\addcontentsline{toc}{subsubsection}{[5] 1.15.2}
  \item[{[5] 1.15.2}] By how much must we improve the CPI of L/S instructions to
        make the program run two times faster?
  \phantomsection\addcontentsline{toc}{subsubsection}{[5] 1.15.3}
  \item[{[5] 1.15.3}] By how much is the execution time of the program improved if
        the CPI of INT and FP instructions is reduced by 40\% and the CPI of
        L/S and Branch is reduced by 30\%?
\end{enumerate}

\problembreak % [5] ≥ 5

%-----------------------------------------------------------------------------
\subsection{[5] Parallel Overhead and Speedup}
\hypertarget{prob-1.16}{}
\noindent$\langle$\S1.8$\rangle$\medskip

\noindent When a program is adapted to run on multiple processors in a multiprocessor
system, the execution time on each processor is comprised of computing time
and the overhead time required for locked critical sections and/or to send
data from one processor to another.

Assume a program requires $t = 100$ s of execution time on one processor.
When run on $p$ processors, each processor requires $t/p$ s, as well as an
additional 4 s of overhead, irrespective of the number of processors. Compute
the per-processor execution time for 2, 4, 8, 16, 32, 64, and 128
processors. For each case, list the corresponding speedup relative to a single
processor and the ratio between actual speedup versus ideal speedup (speedup
if there was no overhead).
